\section{Related Work}\label{sec:related}

\subsection{Chain-of-Thought Distillation to Small Students}

Distillation trains a small student on the outputs of a strong teacher \citep{hinton2015distilling}. Since reasoning models appeared, their long solutions have become the main source of supervision, and the small models distilled and released with R1 have become the starting point for follow-up work \citep{guo2025deepseekr1}.

Reports have accumulated, however, that long solutions can instead hurt small students. Small students learn better from a strong teacher's short solutions than from its long ones \citep{li2025small}, and training on even a small amount of long solutions sharply degrades performance \citep{luo2025valley}. Pairing many teachers with many students shows that the strongest teacher is not always the best, so the cause does not lie with the student alone \citep{hu2026strongest}.

The field has therefore moved from choosing either long or short solutions toward filling the gap between them. Some work adjusts the composition and order of the generated data, as in curricula that mix short and long solutions or change the order of difficulty \citep{jiang2025teach}. Other work turns only the reasoning mode of a single teacher on and off to produce two datasets from the same prompts \citep{boizard2026scale}. When the student is distilled from an interpolated teacher, which mixes a reasoning model and an instruct model in weight space with ratio $\lambda$, it performs best at an intermediate $\lambda$ rather than at either end; the student's solution length changes smoothly with $\lambda$, and long solutions train the student along more concentrated directions than short ones \citep{lan2026midistillation}. The interpolated teacher itself is reported to change its solution style continuously with $\lambda$ until, at some $\lambda$, it stops thinking altogether \citep{wu2026interpolation,lan2025spectrum}. Abrupt emergence of some behaviors depending on the mixing ratio has also been reported when assembling mixture-of-experts models \citep{klagges2025assembly}, while other work reports no such switch \citep{yang2026system}. Reasoning- and instruct-model checkpoints have also been refined jointly through LoRA distillation and weight interpolation, but only accuracy is examined \citep{saini2026thinkswitch}. This line of work, however, selects the teacher by the student's accuracy or evaluates the teacher only by its own outputs, and does not address what remains in the student's weights when the teacher stops thinking.

\subsection{Reading and Editing Model Internals}

Treating a LoRA update as a subspace and comparing different training runs by the overlap of their subspaces dates back to the original LoRA paper \citep{hu2021lora}, and work on SAEs reports that individual directions fluctuate across training runs while the subspace is reproduced \citep{gerasimov2026unstable}.

A line of work opens up trained weights to see what they contain. Training only a rank-1 LoRA on R1's solutions recovers 73--90\% of the performance gap to full fine-tuning, and the adapter's direction is interpreted by reading it through the activations \citep{ward2025rank1}. The teacher's reasoning trajectories themselves are also reported to lie close to a low-dimensional subspace \citep{solgi2026lori}. Spectrally decomposing the weight difference between a reasoning model and an instruct model shows that reasoning ability is concentrated in components lying outside the instruct model's principal singular directions \citep{ma2026spectral}. Between a teacher and a student that share an initialization, the properties transferred by distillation open and close with the alignment of their internal representations \citep{chauhan2026covert}. The subspaces of adapters trained on different tasks overlap substantially, so that an earlier adapter can stand in for a later task, although this overlap is not compared against chance level \citep{dieudonne2026adapter}.

Work that edits the directions it finds to change behavior sometimes treats a direction as a switch that turns the behavior on and off. Removing a single direction that mediates refusal leaves the model unable to refuse, and adding it makes the model refuse even harmless requests \citep{arditi2024refusal}; the same projection removal can also precisely erase a specific capability \citep{fu2026scalpel}. For thinking as well, approximating at low rank what a reasoning model changed from its original model and removing the top component of that approximation reduces overthinking \citep[Gold-Switch;][]{lee2025goldswitch}, and a single activation direction that separates thinking and non-thinking modes is used to move between the two and to measure its effect \citep{liu2026leakage}.

Other work treats the amount of thinking as a dial set by the strength of a direction. One approach trains the magnitude of a single LoRA direction to control reasoning length \citep{ma2025cotvalve}, and another removes a short-reasoning direction in a few attention heads to lengthen reasoning \citep{sun2025thinkedit}. Amplifying reflection SAE features or activation directions lengthens solutions \citep{galichin2026covered,sheng2025planning}, and the strength of a single direction can also set how long the model thinks \citep{jayabahu2026halt}.

Still other work reads internal signals with SAEs to find how a model thinks. Reflection and verification features present only in distilled models are used to steer the model toward overthinking or toward concise thinking \citep{baek2025distilled}. Approximating the difference between a distilled reasoning model and its original model with interpretable components shows that a small number of features producing hesitation markers are necessary and sufficient for those markers, and that removing them shortens responses \citep{hu2026transcoder}. Viewed the same way, RL teaches when to use reasoning mechanisms that already exist, whereas distillation (SFT) implants new ones \citep{venhoff2026base}. Other work counters, however, that many of the reasoning features selected by contrasting long and short solutions respond to the reflective words themselves rather than to thinking \citep{ma2026sae}.

These studies, however, work in settings where the teacher is fixed to a single model or absent altogether (even the studies of distilled models compare one distilled model with its original model), and they find directions by the magnitude of the update or by contrasting long and short solutions. They also treat the directions they find as either a switch or a dial, and do not ask whether the same region responds differently when it is removed, grafted, and amplified. This paper compares, side by side, the adapters of students distilled from a teacher moved along $\lambda$ to find the shared region that forms in the student's weights when the teacher stops thinking, and then removes, grafts, or amplifies that region to test its relation to thinking.
